\documentclass[11pt]{article}

\usepackage[margin=1in]{geometry}
\usepackage{amsmath, amssymb}
\usepackage{algorithm}
\usepackage{algpseudocode}
\usepackage{booktabs}
\usepackage{hyperref}
\usepackage{braket}

\title{Intro to Quantum Computing - Week 1 Notes}
\author{}
\date{}

\begin{document}
\maketitle

\section{Quantum Postulates}

\begin{itemize}
    \item Postulate one: How to model states: Every closed quantum state can be represented by a vector $\ket{\psi}$ in a vector space in $ \mathbb{C}^n $ ($\mathbb{R}^n$), this can be $\ket{\psi}=\begin{bmatrix} c_1 \\ c_2 \\ \vdots \\ c_n \end{bmatrix}$, where the vectors are noralized, meaning that $|c_1|^2 + |c_2|^2 + ... + |c_n|^2 = 1$. If we have a collection of vectors whch are part of the set of V, that is $\ket{r} + \ket{w} \in V$, for $\alpha\ket{v} \in V$
    \item Postulate two: How can we measure: Measurements are described by a set of operators that act on the state vector.
    \item Postulate three: State evolution: The time evolution of a closed quantum system is governed by the Schrödinger equation.
    \item Postulate four: Composite systems: The state of a composite quantum system is the tensor product of the states of its components.
\end{itemize}


cLosed systemes are deterministic
\end{document}
